\teachcover{Cyclic groups\\and counting matrices}{1}{
Multiplication in a finite field has a generator
\[
K^\times=\{1,g,g^2,\ldots,g^{q-2}\},\qquad |K|=q.
\]
}{
This lecture proves cyclicity of a finite field's multiplicative group
and then counts ordered bases and determinant fibers. Recall Lagrange's
theorem and the definition of element order. We prove the polynomial
root bound before using it. The group-size-one case is included.
The optional slides compare finite fields with units of prime-power
residue rings; they are not part of the 50-minute main route.
}

\teachslide{The nonzero elements form a group}{2}{
Let $K$ be a field with $q$ elements.
\[
K^\times=K\setminus\{0\},\qquad |K^\times|=N=q-1.
\]
Multiplication is commutative; the identity is $1$;
every element has an inverse.
\[
\key{\ord(a)\mid N,\qquad a^N=1\quad(a\in K^\times).}
\]
}{
Closure follows from absence of zero divisors: the product of two
nonzero elements cannot be zero. Inverses are nonzero. Lagrange
gives the displayed order divisibility. Define order as the least
positive exponent returning to $1$. The task is to find one element
of order $N$, which is stronger than all orders dividing $N$.
}

\teachslide{What cyclic means}{2}{
In $\F_5$, the powers of $2$ are
\[
\begin{array}{c|rrrr}
t&0&1&2&3\\\hline
2^t&1&2&4&3
\end{array}
\qquad 2^4=1.
\]
Every nonzero element is a power of $2$.
\spacer
A group is \key{cyclic} if one element generates the whole group.
}{
Calculate the last power by reducing $8$ modulo $5$. A generator of
a finite group of size $N$ has order $N$, and conversely an element
of order $N$ supplies all $N$ elements by its distinct powers.
There may be several generators. Our aim is to prove at least one
exists in every finite field, not to identify a canonical generator.
}

\teachslide{A polynomial cannot have too many roots}{2}{
Over a field, if $f(a)=0$, then
\[
f(X)=(X-a)h(X).
\]
Any other root $b\ne a$ is a root of $h$:
\[
0=(b-a)h(b)\quad\Longrightarrow\quad h(b)=0.
\]
Induction gives:
\[
\key{\deg f=d,\ f\ne0\quad\Longrightarrow\quad\text{at most }d\text{ roots}.}
\]
}{
Recall polynomial division by $X-a$. The constant nonzero case has
no roots. Each root factors off one linear term; cancellation works
since $b-a$ is a nonzero field element. This proof is the specific
place field arithmetic enters the cyclicity theorem. The root bound
can fail over a ring with zero divisors.
}

\teachslide{Build the order one prime at a time}{2}{
Factor $N=q-1$:
\[
N=\ell_1^{r_1}\cdots\ell_t^{r_t}.
\]
We will construct, for each prime factor $\ell$,
\[
\key{g_\ell\text{ with order }\ell^r,\qquad \ell^r\parallel N.}
\]
Then multiply these elements.

If $N=1$, the group $\{1\}$ is already cyclic.
}{
Explain the double-bar notation: $\ell^r$ divides $N$, but the next
power does not. Distinct primes have relatively prime powers. The
product step requires both coprime orders and commutativity, proved
below. For $K=\F_2$, $N=1$ and $1$ is a generator of the trivial group.
}

\teachslide{Find an element with the required prime factor}{3}{
Fix $\ell^r\parallel N$.

The polynomial
\[
X^{N/\ell}-1
\]
has at most $N/\ell<N$ roots.

It cannot vanish at every element of $K^\times$.
\[
\key{\text{Choose }a\in K^\times\text{ with }a^{N/\ell}\ne1.}
\]
}{
Do not assume a generator to choose $a$: the root bound proves that
at least one of the $N$ nonzero elements fails this equation. The
chosen $a$ can depend on $\ell$, which is fine. Its order must include
the full $\ell$-power of $N$, but the next slide proves the needed
statement directly through powers rather than invoking an order formula.
}

\teachslide{Isolate the prime-power order}{3}{
Set $M=N/\ell^r$ and $g_\ell=a^M$.
\[
g_\ell^{\ell^r}=a^N=1,
\qquad g_\ell^{\ell^{r-1}}=a^{N/\ell}\ne1.
\]
Its order divides $\ell^r$, but does not divide $\ell^{r-1}$.

Every proper divisor of $\ell^r$ divides $\ell^{r-1}$.
\[
\key{\ord(g_\ell)=\ell^r.}
\]
}{
Read both exponent products aloud. If $x^m=1$, its order divides
$m$: divide $m$ by the least positive return time and use minimality
of the remainder. This fact is more precise than Lagrange and also
applies here to a known exponent. The prime-power divisor observation
is what forces the exact order without computing it by trial.
}

\teachslide{Combine two coprime orders}{4}{
Suppose $a,b$ commute, with orders $r,s$ and $\gcd(r,s)=1$.
\[
(ab)^{rs}=1.
\]
If $(ab)^t=1$, put $z=a^t=b^{-t}$. Then
\[
z^r=z^s=1\quad\Longrightarrow\quad z=1
\]
by Bezout's identity. Thus $r\mid t$ and $s\mid t$, so $rs\mid t$.
\[
\key{\ord(ab)=rs.}
\]
}{
This lemma deserves the full four minutes. If $ur+vs=1$, then
$z=z^{ur+vs}=(z^r)^u(z^s)^v=1$; negative exponents use inverses.
Hence $a^t=b^t=1$. The order-divides-return-time fact forces both
divisibilities, whose coprimality gives $rs\mid t$. We already know
$rs$ is a return time, so it is the least. Commutativity is needed
for $(ab)^t=a^tb^t$. The result does not claim that arbitrary
products have the lcm of their factors' orders.
}

\teachslide{One element of full order}{3}{
The prime-power orders $\ell_i^{r_i}$ are pairwise coprime.

By repeatedly applying the lemma,
\[
g=g_{\ell_1}\cdots g_{\ell_t}
\quad\text{has order}\quad
\ell_1^{r_1}\cdots\ell_t^{r_t}=N.
\]
Its $N$ distinct powers exhaust $K^\times$.
\[
\key{\text{The multiplicative group of every finite field is cyclic.}}
\]
}{
Check that a partial product's order remains coprime to the next
prime-power order, so induction applies. An element of order $N$
has exactly $N$ distinct powers from exponent $0$ through $N-1$;
there are only $N$ elements in the group. The proof has used Lagrange,
the polynomial root bound, Bezout, and commutativity, with no general
classification theorem for finite abelian groups.
}

\teachslide{Primitive elements}{2}{
A generator $g$ of $K^\times$ is called a \key{primitive element}.
\[
K^\times=\{1,g,\ldots,g^{q-2}\},\qquad \ord(g)=q-1.
\]
If $K$ has characteristic $p$, then
\[
\key{K=\F_p(g):\text{ the field is generated by one element}.}
\]
Every nonzero element is a power of $g$.
}{
In this finite-field context primitive means generator of the
multiplicative group. Such an element also generates the extension
field, since a field containing it contains all its powers and zero.
The converse need not hold: generating a field extension does not
necessarily mean having multiplicative order $q-1$. In particular,
the Fibonacci root modulo $7$ will have order $16$, not $48$.
}

\teachslide{Your turn: combine orders in $\F_7$}{3}{
In $\F_7^\times$,
\[
\ord(6)=2,\qquad\ord(4)=3.
\]
Verify these orders, then answer:
\begin{enumerate}
\item What is the order of $6\cdot4$?
\item Which residue is therefore a primitive element?
\end{enumerate}
}{
Give a minute of calculation. Students should find $6^2=1$ with
$6\ne1$, and $4^2=2$, $4^3=1$ with $4\ne1$. Coprimality lets us
combine the orders. This is a concrete version of the theorem's
final construction, not a search through all six possible generators.
}

\teachslide{A generator modulo seven}{1}{
\[
6\cdot4=3\pmod7,
\qquad \ord(3)=2\cdot3=6.
\]
\[
\begin{array}{c|rrrrrr}
t&0&1&2&3&4&5\\\hline
3^t&1&3&2&6&4&5
\end{array}
\]
\key{One element generates all six nonzero residues.}
}{
The table verifies the conclusion independently. We used the lemma
to know the order before listing all powers. This ends the cyclicity
portion at minute 28. The rest of this lecture counts matrix groups.
}


\teachslide{From field elements to matrices}{2}{
Over $\F_p$, the span of $r$ independent vectors has $p^r$ elements.
\spacer
An invertible matrix has columns forming an ordered basis.
\[
\key{\text{Count invertible matrices by choosing their columns.}}
\]
}{
Retrieve the coordinate count from Lecture 2. The remaining 20 minutes
count GL and SL, providing the group-size bounds we will use for the
Fibonacci matrix in Lecture 4. The cyclicity theorem concerns field
elements under multiplication, not a claim that GL or SL is cyclic.
}

\teachslide{Invertible matrices and ordered bases}{2}{
For a matrix with columns $v_1,\ldots,v_d$,
\[
M\begin{pmatrix}a_1\\\vdots\\a_d\end{pmatrix}
=a_1v_1+\cdots+a_dv_d.
\]
$M$ is invertible exactly when every output has unique coordinates.
\spacer
\key{Its columns must form an ordered basis of $\F_p^d$.}
}{
Connect matrix invertibility to the bijectivity of its linear map.
A basis supplies existence and uniqueness of solutions to $Mx=y$.
Conversely, an invertible map makes the columns spanning and independent.
Over a field this also agrees with nonzero determinant. Order matters:
exchanging columns usually gives another matrix.
}

\teachslide{Choose the first column}{2}{
In $\F_p^2$, there are $p^2$ vectors.

The first basis vector can be anything except zero.
\[
\key{p^2-1\text{ choices}.}
\]
\[
v_1\ne\begin{pmatrix}0\\0\end{pmatrix}.
\]
}{
Ask students why a zero column prevents invertibility. A singleton
list containing a nonzero vector is independent because a nonzero
scalar has an inverse. This will matter when we count its span.
}

\teachslide{Choose the second column}{3}{
Once $v_1\ne0$ is fixed, its span has $p$ elements:
\[
\Span(v_1)=\{av_1:a\in\F_p\}.
\]
Choose $v_2$ outside that span:
\[
\key{p^2-p\text{ choices}.}
\]
Therefore
\[
|\GL_2(\F_p)|=(p^2-1)(p^2-p).
\]
}{
The scalar multiples are distinct because $v_1$ is nonzero. Being
outside the span gives independence, by the basis-extension argument
from Lecture 2. Two independent vectors in $\F_p^2$ span it: their
span has $p^2$ elements, the full ambient size. The number of allowed
second columns does not depend on which first column was chosen.
}

\teachslide{Your turn: count modulo two}{2}{
How many invertible $2\times2$ matrices are there over $\F_2$?

Fix the first column to be
\[
v_1=\begin{pmatrix}1\\0\end{pmatrix}.
\]
Which second columns are allowed?
}{
Give quiet work time, then ask for both the total count and the actual
two allowed columns for this fixed choice. Students should list the
four vectors and exclude $0$ and $v_1$. This reinforces the reason
for subtraction in the count.
}

\teachslide{Six invertible matrices}{1}{
\[
|\GL_2(\F_2)|=(4-1)(4-2)=6.
\]
With $v_1=(1,0)^{\mathsf T}$, the allowed second columns are
\[
\begin{pmatrix}0\\1\end{pmatrix},\qquad
\begin{pmatrix}1\\1\end{pmatrix}.
\]
The span we excluded is $\{(0,0),(1,0)\}$.
}{
Confirm that zero and the first column are exactly the forbidden
vectors. Over $\F_2$, every nonzero determinant is $1$, so GL and
SL will have the same size. Mention this only as an anticipation.
}

\teachslide{The same count in dimension $d$}{2}{
After $r$ independent columns, the span has $p^r$ vectors.

The next column has $p^d-p^r$ choices.
\[
\key{|\GL_d(\F_p)|=\prod_{r=0}^{d-1}(p^d-p^r).}
\]
The first factor is $p^d-1$; the last is $p^d-p^{d-1}$.
}{
The span of the empty list has size $1$, accounting for the first
factor when $r=0$. Multiplying the independent choices counts ordered
bases. Do not expand or simplify the product: the unsimplified form
records the argument. Most lecture time stays with $d=2$.
}

\teachslide{The determinant groups matrices by value}{2}{
\[
\det:\GL_d(\F_p)\longrightarrow\F_p^\times.
\]
\[
\det(MN)=\det(M)\det(N).
\]
Every nonzero $a\in\F_p$ occurs:
\[
D_a=\operatorname{diag}(a,1,\ldots,1),\qquad\det D_a=a.
\]
\key{$\SL_d(\F_p)$ is the set with determinant $1$.}
}{
Define $\F_p^\times$ as the multiplicative group of nonzero field
elements, of size $p-1$. The determinant is a surjective group
homomorphism. Quotient-group notation is optional; the next slide
uses a concrete bijection between determinant fibers.
}

\teachslide{Every determinant fiber has the same size}{2}{
Multiplication by $D_a$ gives a bijection
\[
\SL_d(\F_p)\longrightarrow\{M:\det M=a\},
\qquad S\longmapsto D_aS.
\]
The inverse is $M\longmapsto D_a^{-1}M$.

There are $p-1$ nonzero determinant values, so
\[
\key{|\SL_d(\F_p)|=\frac{|\GL_d(\F_p)|}{p-1}.}
\]
}{
Verify the two determinants: $\det(D_aS)=a$ and
$\det(D_a^{-1}M)=a^{-1}a=1$. Fibers are disjoint and exhaust GL.
This establishes the count without assuming the first isomorphism
theorem. If that theorem is familiar, the quotient is on a backup slide.
}

\teachslide{In dimension two}{2}{
\[
\begin{aligned}
|\SL_2(\F_p)|
&=\frac{(p^2-1)(p^2-p)}{p-1}\\
&=\key{p(p^2-1)}.
\end{aligned}
\]
Lagrange's theorem: the order of a group element divides the group size.
\spacer
\key{Next time: apply this count to Fibonacci.}
}{
Factor $p^2-p=p(p-1)$ before cancelling. Briefly recall the element-order
definition: the least positive exponent giving the identity. Lagrange
applies to the cyclic subgroup generated by an element. This assumes
the course's earlier group material rather than reproving Lagrange.
}

\teachslide{Exit question}{2}{
For a field $K$ with $q$ elements, which group have we proved is cyclic?

What is its order?

Why does this not say that every nonzero element generates it?
}{
Expected answers: $K^\times$, with multiplication, of order $q-1$.
The proof finds an element of full order; many other elements can have
smaller orders, including $1$ when $q>2$. Ask students also to recall
the size $p(p^2-1)$ of SL(2,p). End at minute 50.
}

\backups

\backupslide{The determinant quotient}{
The determinant is surjective, with kernel $\SL_d(\F_p)$.
\[
\key{\GL_d(\F_p)/\SL_d(\F_p)\cong\F_p^\times.}
\]
A coset is determined by the common determinant of its matrices.

Multiplication of cosets corresponds to multiplication of determinants.
}{
This is the first isomorphism theorem. Kernel subgroups are normal,
so the quotient is defined. The explicit fiber bijection in the main
deck already proved the group-size formula without needing quotient
terminology. This optional slide answers the literal GL/SL question.
}
\backupslide{A different object with a cyclic unit group}{
For odd primes $p$ and $n\ge1$,
\[
\key{(\Z/p^n\Z)^\times\text{ is cyclic of order }p^{n-1}(p-1).}
\]
For $n>1$, the ring is not a field: $p\cdot p^{n-1}=0$.
\[
\begin{array}{c|c|c}
&\F_{p^n}&\Z/p^n\Z\\\hline
\text{Units}&p^n-1&p^{n-1}(p-1)\\
\text{Characteristic}&p&p^n
\end{array}
\]
}{
This is additional material, not an application of the finite-field
theorem to a residue ring. Units are exactly the classes not divisible
by $p$, giving $p^n-p^{n-1}$. Both factors in the displayed zero product
are nonzero when $n>1$. The next slides give a primitive-root lifting
proof. The field comparison assumes a field of order $p^n$; we have
not proved existence for every prime power in the main lectures.
}

\backupslide{Two arithmetics with nine elements}{
In $\Z/9\Z$, the powers of $2$ are
\[
1,2,4,8,7,5,\qquad2^6=1.
\]
These are all six units; the nonzero residues $3,6$ are not units.
\spacer
In $\F_9$, all eight nonzero elements are units,
and a primitive element has order $8$.
}{
Calculate successive powers modulo $9$. The nonzero residue $3$
satisfies $3^2=0$. A field of nine elements can instead be constructed
by adjoining a square root of $2$ to $\F_3$, using Lecture 2's method.
This contrasts a six-element unit group with an eight-element
multiplicative field group, despite equal sizes of the underlying sets.
}

\backupslide{First lift a primitive root to modulo $p^2$}{
Let $p$ be odd, and let $g$ generate $\F_p^\times$.

Choose $h=g$ or $h=g+p$ so that
\[
\key{h^{p-1}\not\equiv1\pmod{p^2}.}
\]
This is possible because
\[
(g+p)^{p-1}\equiv g^{p-1}+p(p-1)g^{p-2}\pmod{p^2},
\]
and the added term is nonzero modulo $p^2$.
}{
Choose any integer representative $g$ prime to $p$. The two choices
have the same primitive reduction modulo $p$. The linear binomial
term is divisible by $p$ but not $p^2$; all remaining terms are
divisible by $p^2$. Thus both powers cannot equal 1 modulo $p^2$.
Write $h^{p-1}=1+pc$ with $p\nmid c$. It would be false to say
every primitive representative automatically lifts without adjustment.
}

\backupslide{Each further lift contributes a factor of $p$}{
For odd $p$, if $u=1+p^rc$ with $r\ge1$ and $p\nmid c$, then
\[
u^p=1+p^{r+1}c',\qquad p\nmid c'.
\]
The first binomial term is $p^{r+1}c$; all later terms
are divisible by $p^{r+2}$.
\spacer
Starting with $h^{p-1}=1+pc$ gives
\[
\key{v_p\bigl(h^{(p-1)p^k}-1\bigr)=k+1.}
\]
}{
Define $v_p(m)$ as the exponent of $p$ dividing the nonzero integer
$m$, if needed. For $2\le j<p$, the binomial coefficient contributes
$p$, so the $j$th term has valuation at least $1+rj\ge r+2$.
The last term has valuation $rp\ge r+2$ since $p\ge3$.
After division by $p^{r+1}$, the result is $c$ modulo $p$, so
it remains nonzero. Induction gives the displayed formula.
This is a direct special-case proof, not an invocation of a general
lifting-the-exponent lemma.
}

\backupslide{The lifted element has full order}{
Modulo $p^n$, let $d=\ord(h)$.
\[
p-1\mid d,\qquad d\mid(p-1)p^{n-1}.
\]
Thus $d=(p-1)p^j$ for some $0\le j\le n-1$.

The lifting calculation shows the least such $j$ is $n-1$.
\[
\key{d=(p-1)p^{n-1}.}
\]
So $h$ generates every unit modulo $p^n$.
}{
Reduction modulo $p$ forces $p-1$ to divide the order; Lagrange
in the unit group gives the upper divisibility. The valuation formula
says $h^{(p-1)p^j}\equiv1\pmod{p^n}$ exactly when $j+1\ge n$.
This proves the exact order, with a single adjusted integer $h$
working for every $n$. The case $n=1$ is the finite-field theorem.
}

\backupslide{The exception starts at eight}{
\[
(\Z/2\Z)^\times=\{1\},\qquad(\Z/4\Z)^\times=\{1,3\}.
\]
Both are cyclic. Modulo $8$, every unit squares to $1$:
\[
1^2=3^2=5^2=7^2=1\pmod8.
\]
There is no element of order $4$.

For $n\ge3$,
\[
\key{(\Z/2^n\Z)^\times\cong C_2\times C_{2^{n-2}}.}
\]
}{
$C_m$ denotes a cyclic group of order $m$. The direct-product
generators are $-1$ and $5$. Repeated squaring gives
$v_2(5^{2^k}-1)=k+2$, starting with $5-1=4$:
if $u=1+2^rc$, $r\ge2$, $c$ odd, then
$u^2-1=2^{r+1}c(1+2^{r-1}c)$ and the last factor is odd.
So 5 has order $2^{n-2}$, and its powers are exactly the residues
equal to 1 modulo 4. Every odd residue is in this subgroup or its
negative, and $-1$ is outside it. This proves the direct product.
Its largest element order is $2^{n-2}$, smaller than its size
$2^{n-1}$, so it is not cyclic. The modulo-8 example alone suffices
for a short classroom comparison.
}

\backupslide{Why the field proof does not apply to residue rings}{
In $\Z/8\Z$, the degree-two polynomial
\[
X^2-1
\]
has four roots: $1,3,5,7$.
\spacer
\key{The polynomial root bound required a field.}

An abelian unit group need not be cyclic.
}{
The finite-field proof chose an element failing a low-degree polynomial
equation. With zero divisors, cancellation fails in the factor-theorem
argument. Thus the residue-ring theorem needs an additional lifting
argument, and the theorem fails for powers of 2 starting at 8.
This returns to Lecture 1's discussion of zero divisors.
}


